\chapter{Introduction}
\label{ch:introduction}

\section{Overview}

Work no longer happens in one place. Over the last decade, and dramatically so
after 2020, teams have moved from a shared physical office to a distributed
arrangement in which colleagues may be in different homes, cities, or time
zones. The tools that grew up to support this transition are powerful but
\emph{disjoint}: a team may hold meetings in one application, chat in a second,
track tasks in a third, keep its calendar in a fourth, and store files in a
fifth. Each tool is a silo with its own identity, its own notifications, and its
own mental model. The result is a workday spent switching between browser tabs
and losing the two things a real office provides for free: an ambient sense of
\emph{who is present and available}, and the \emph{spontaneous} conversations
that begin simply because two people happen to be near each other.

\voname{} is a graduation project that responds to this problem with a single
idea: give a distributed team a \emph{place} rather than a collection of tools.
That place is a two-dimensional, game-like office --- a digital headquarters ---
in which every member is represented by an animated avatar. Users walk around a
tile-mapped floor, see their colleagues moving in real time, sit down at desks,
gather around whiteboards, and slip into meeting corners. When two avatars move
close together, a proximity-based voice and video call opens automatically, just
as a hallway conversation would; when they move apart, it ends. Around this
spatial core, \voname{} folds in the rest of the collaboration stack ---
persistent chat channels and direct messages, threaded discussions, task
assignment, calendar scheduling, room-based meetings, and a unified notification
system --- so that the entire working day can happen inside one product.

Delivering this experience is a substantial distributed-systems undertaking. The
platform comprises an animated multiplayer world, a full-featured web client, a
native-quality mobile client, and a backend of nine cooperating microservices,
all bound together by real-time messaging over WebSockets, event-driven
communication over a message bus, and a shared identity model enforced at the
gateway. This book is a complete engineering account of that system.

\section{Motivation}

The motivation for \voname{} rests on three observations about modern
distributed work.

\begin{description}
  \item[Tool fragmentation is expensive.] The average knowledge worker
  now moves between many single-purpose applications during a day. Each context
  switch carries a cognitive cost, and each additional tool multiplies the
  surfaces on which information can be lost. Consolidating presence,
  communication, tasks, and scheduling into one platform removes friction and
  gives a team a single source of truth for ``what is happening right now.''

  \item[Video calls kill spontaneity.] Scheduled video meetings are excellent
  for planned discussion but hostile to the unplanned. In a physical office,
  most valuable communication is incidental --- a question asked in passing, an
  overheard problem, a quick huddle at a desk. Grid-of-faces video conferencing
  has no equivalent: you cannot ``walk past'' someone in a Zoom call. A spatial
  interface, where being near a colleague is enough to start talking, restores
  this incidental channel.

  \item[Presence is a first-class need.] Knowing who is online, who is busy, and
  who is deep in focus is information a physical office broadcasts continuously
  and effortlessly. Most remote tools reduce this to a coloured dot. A spatial
  office makes presence \emph{legible}: you can literally see where people are,
  whether they are at their desk, in a meeting room, or gathered with others.
\end{description}

Commercial products such as Gather and other ``metaverse office'' tools have
validated that a spatial approach resonates with users. \voname{} is motivated
by the goal of building such a platform from first principles --- not as a thin
demo, but as a properly architected, multi-client, microservices-based system ---
and, in doing so, of learning and demonstrating the engineering that a real
real-time collaboration product requires.

\section{Problem Statement}

The project sets out to answer a single engineering question:

\begin{infobox}[title=Problem Statement]
How can a distributed team be given a unified digital workplace that reproduces
the \emph{presence} and \emph{spontaneity} of a physical office, while
consolidating the fragmented tools of remote work --- messaging, meetings,
tasks, scheduling, and workspace administration --- into one coherent,
real-time, multi-client platform?
\end{infobox}

Answering it well requires solving several sub-problems that recur throughout
this book:

\begin{enumerate}
  \item \textbf{Real-time spatial state.} Many users must see each other move
  around a shared world with low latency and without the network cost of
  broadcasting every pixel of movement to everyone.

  \item \textbf{Proximity-driven communication.} Voice, video, and screen
  sharing must start and stop based on where avatars are, and must scale without
  routing media through the application's own servers.

  \item \textbf{Reliable messaging.} Chat must deliver instantly, survive
  disconnections, avoid duplicates, and preserve history --- across both a web
  and a mobile client.

  \item \textbf{A coherent identity and permission model} that is enforced once,
  at the edge, and trusted by every internal service.

  \item \textbf{Loose coupling between services} so that features such as
  notifications and room--chat binding can evolve independently and survive
  partial outages.
\end{enumerate}

\section{Objectives}

The concrete objectives of the \voname{} project are:

\begin{enumerate}
  \item To design and implement an \textbf{immersive 2D virtual office} in which
  multiple users, represented as avatars, move and interact in real time within
  a shared, persistent space.

  \item To provide \textbf{proximity-based real-time communication} --- voice,
  video, and screen sharing --- that begins and ends according to spatial
  closeness, using an external media infrastructure so the platform's own
  servers never carry audio or video streams.

  \item To build a \textbf{full-featured messaging system} supporting workspace
  channels, direct messages, threaded conversations, mentions, replies, edits,
  soft deletion, read receipts, and typing indicators, delivered over WebSockets
  with a REST fallback.

  \item To implement a \textbf{workspace and identity layer} that models
  companies, teams, desks (user seats), roles, and invitations, and that issues
  and validates authentication centrally.

  \item To add the surrounding \textbf{collaboration services} --- task
  assignment and tracking, calendar scheduling, meeting rooms, and a unified
  notification centre (both e-mail and in-app).

  \item To expose the platform through \textbf{three clients} --- a 2D
  multiplayer world, a React web application, and a Flutter mobile application ---
  that share the same backend.

  \item To realise all of the above on a \textbf{microservices architecture}
  that demonstrates sound distributed-systems practice: an API gateway, per-service
  data ownership, event-driven decoupling, and consistent cross-cutting patterns
  for security and real-time communication.
\end{enumerate}

\section{Scope of the Project}

\voname{} is a large system, and it is important to state clearly what this book
documents and where the boundaries lie.

\paragraph{In scope.} The book covers the full breadth of the delivered system:
the nine backend microservices and their data models, APIs, and message flows;
the 2D multiplayer office (client and server); the React web client; and the
Flutter mobile client. It covers the cross-cutting mechanisms that make the
platform work as a whole --- gateway authentication, token-based WebSocket
handshakes, the event bus, and the real-time messaging protocol.

\paragraph{Boundaries and honest limitations.} \voname{} is a graduation
project, and several parts of the system are at differing levels of maturity. As
this book documents faithfully, some capabilities are implemented end-to-end
(chat, workspaces, the 2D office, rooms), some are implemented on one client but
not another (the mobile app implements only the chat domain), some backend
services are more complete than others, and some features are deliberately left
as ``version~2'' hardening (for example, server-side minting of media tokens, or
role-based authorisation checks that are currently trusted to the client).
Rather than gloss over these, the book records them as they are; they are
valuable both as an accurate description of the system and as a guide to future
work.

\section{Main Contributions}

The principal contributions of the project, and of this book, are:

\begin{itemize}
  \item A complete, working design for a \textbf{spatial collaboration platform}
  that unifies a real-time 2D office with a full messaging, tasks, calendar, and
  workspace stack.

  \item A \textbf{nine-service microservices backend} built on Spring~Boot,
  demonstrating an API gateway with centralised authentication, polyglot
  persistence (PostgreSQL, MongoDB, Redis), and event-driven decoupling over
  RabbitMQ.

  \item A set of \textbf{reusable real-time patterns} --- ticket-based WebSocket
  authentication, a uniform \code{\{action, payload\}} broadcast envelope, and
  proximity-triggered peer-to-peer media --- applied consistently across
  services and clients.

  \item Three \textbf{distinct client implementations} (2D game, web, mobile)
  over one backend, illustrating how a single platform can be surfaced through
  very different technologies.

  \item A thorough \textbf{engineering write-up} that can serve as a reference
  for building real-time, multi-client collaboration systems.
\end{itemize}

\section{Organisation of the Book}

The remainder of this book is organised as follows.

\begin{description}
  \item[Chapter~\ref{ch:background}] surveys the background and related work: the
  concept of the virtual office, the enabling technologies (WebSockets, WebRTC,
  microservices, game engines), and a comparison with existing tools.

  \item[Chapter~\ref{ch:requirements}] presents the requirements analysis ---
  actors, functional and non-functional requirements, user stories, and use-case
  and domain models expressed in UML.

  \item[Chapter~\ref{ch:architecture}] describes the overall system
  architecture: the microservices, the gateway, the data and messaging
  infrastructure, and the cross-cutting patterns that unify them.

  \item[Chapter~\ref{ch:backend}] documents each backend microservice in
  depth --- its responsibility, data model, REST and WebSocket API, and its
  event interactions with the rest of the system.

  \item[Chapter~\ref{ch:skyoffice}] details the 2D virtual office: its
  client/server split, the multiplayer synchronisation model, interactive
  objects, and the proximity-based communication stack.

  \item[Chapter~\ref{ch:web}] and Chapter~\ref{ch:mobile} cover the web and
  mobile clients respectively, including their architecture, real-time
  integration, and feature sets.

  \item[Chapter~\ref{ch:testing}] discusses the testing and evaluation
  approach, and Chapter~\ref{ch:conclusion} concludes with an assessment of the
  project and a roadmap for future work.
\end{description}
